\documentclass[11pt]{article}
\usepackage[a4paper, margin=1in]{geometry}
\usepackage{amsmath, amssymb, bm}
\usepackage{hyperref}

\title{Stage-5 Deployment Math:\\
       Latency Budget, Success-Rate Translation, and the Stateless-vs-Stateful Choice}
\author{REACT stage-5 pre-flight}
\date{2026-05-22}

\newcommand{\R}{\mathbb{R}}
\newcommand{\ms}{\,\text{ms}}
\newcommand{\flops}{\,\text{FLOPS}}

\begin{document}
\maketitle

\section{The two paper-target numbers}

REACT's paper-level targets (set in README, traced from PEMTRS RA-L 2026):
\begin{enumerate}
  \item Inference latency $T_{\text{inf}} < 10\ms$ per frame on a Jetson-class
        device (Orin Nano or NX as the typical reference).
  \item Closed-loop success rate $\text{SR} \geq 85\%$ in dynamic-obstacle
        scenes.
\end{enumerate}

The current training metric \texttt{dyn\_dyn\_tail} $\approx 0.22$--$0.24$ is
\emph{not} either of these. This note derives the bridge.

\section{Latency budget decomposition}

Per-frame inference, in temporal mode (K=10 stateless):
\begin{equation}\label{eq:latency-stateless}
T_{\text{inf}}^{\text{stateless}}
\;=\;
T_{\text{ReVAE-enc}}^{K} + T_{\text{GRU}}^{K} + T_{\text{backbone}}^{1} +
T_{\text{head}}^{1} + T_{\text{post}}.
\end{equation}

Single-frame mode (stage-3.1):
\begin{equation}\label{eq:latency-single}
T_{\text{inf}}^{\text{single}}
\;=\;
T_{\text{ReVAE-enc}}^{1} + T_{\text{backbone}}^{1} +
T_{\text{head}}^{1} + T_{\text{post}}.
\end{equation}

Stateful mode (DiffPhysDrone-style, see \texttt{stage\_3\_4\_research\_cn.md}):
\begin{equation}\label{eq:latency-stateful}
T_{\text{inf}}^{\text{stateful}}
\;=\;
T_{\text{ReVAE-enc}}^{1} + T_{\text{GRU}}^{1} + T_{\text{backbone}}^{1} +
T_{\text{head}}^{1} + T_{\text{post}}.
\end{equation}

\subsection{Per-component estimates (Orin NX, INT8 quantised reVAE)}

Based on
\href{https://developer.nvidia.com/embedded/jetson-orin}{Jetson Orin NX} INT8
throughput (TF32 baseline $\approx 70$ TFLOPS at FP16) and the module FLOP
counts (\texttt{reVAE.encode} $\approx 0.4$ GFLOPS per frame at $96\times160$
input; \texttt{YopoBackbone} ResNet18 $\approx 1.8$ GFLOPS at the same input):
\begin{center}
\begin{tabular}{lrrr}
\hline
Component & FLOPS / call & Time / call (FP16) & Time / call (INT8) \\
\hline
\texttt{reVAE.encode}     & 0.4 G & 0.50\ms & 0.20\ms \\
\texttt{reVAE.decode}     & 0.4 G & not on inference graph & --- \\
\texttt{nn.GRU} step      & 0.1 M & 0.05\ms & 0.05\ms \\
\texttt{YopoBackbone}     & 1.8 G & 2.4\ms  & 1.0\ms \\
\texttt{YopoHead}         & 5.0 M & 0.08\ms & 0.04\ms \\
post-processing (poly solve, anchor argmax) & --- & 1.0\ms & 1.0\ms \\
\hline
\end{tabular}
\end{center}

\subsection{Latency budget table}

Plugging in (FP16):
\begin{center}
\begin{tabular}{lrl}
\hline
Mode                    & Latency & $<10\ms$ budget hit? \\
\hline
\texttt{single} (3.1)   & $0.5 + 2.4 + 0.08 + 1.0 \approx 4.0\ms$  & Yes (with headroom) \\
\texttt{stateful} K=1   & $0.5 + 0.05 + 2.4 + 0.08 + 1.0 \approx 4.0\ms$ & Yes \\
\texttt{stateless} K=10 & $5.0 + 0.5 + 2.4 + 0.08 + 1.0 \approx 9.0\ms$  & Marginal (10\% margin) \\
\hline
\end{tabular}
\end{center}

Conclusion: \textbf{stage-3.4 stateless K=10 is feasible but uncomfortable on
the Orin NX FP16 path}. Either move to INT8 (cuts reVAE encode to 0.2\,ms
each, giving $\sim 6.0\ms$ total at K=10) or switch to stateful at deployment
(stays at $\sim 4.0\ms$). The stateful refactor is \emph{not} a model
re-training --- it's an inference-loop change: feed depth one frame at a time,
carry GRU hidden state between calls.

\section{From \texttt{dyn\_dyn\_tail} to success rate}

This is the underspecified link. We have:
\[
\texttt{dyn\_dyn\_tail} = 0.23 \quad \text{(stage-3.1+, v3 dataset)}.
\]
We want to predict $\text{SR}$ in closed-loop sim. There is no clean closed
form because $\text{SR}$ depends on (i) the closed-loop policy stability,
(ii) the dynamic scene difficulty, and (iii) the controller tracking error
on top of the planner. But we can bound it.

\subsection{Per-trajectory collision probability under a loss-rate prior}

Let $p_c$ denote the per-frame probability that the predicted endstate
$\bm{p}_i^{\text{pred}}$ falls inside a forbidden region (within $d_{\text{safe}}$
of some obstacle). Under the assumption that the optimal anchor is selected
with high probability and \texttt{dyn\_dyn} is roughly the mean penalty over
active anchors:
\begin{equation}\label{eq:pc-estimate}
\texttt{dyn\_dyn}
\;\approx\;
\lambda_{\text{dyn}} \cdot p_c \cdot \overline{\softplus} ,
\end{equation}
where $\overline{\softplus}$ is the mean softplus value on active anchors
$\approx 0.3$. With $\lambda_{\text{dyn}} = 3.0$ and
\texttt{dyn\_dyn} $= 0.23$, this gives
\begin{equation}
p_c \approx \frac{0.23}{3.0 \cdot 0.3} \approx 0.26.
\end{equation}
Interpretation: in $\approx 26\%$ of frames, \emph{some} candidate anchor
emits an endstate violating the safety margin. The selected (top-score)
anchor is usually a safer one, so per-frame \emph{actual} collision
probability is significantly lower.

\subsection{Translating to closed-loop $\text{SR}$}

Assume a closed-loop trajectory consists of $N_{\text{frame}} \approx 100$ planning
steps (3.3\,s flight at 30\,Hz), and that the score head selects the
\emph{best} anchor each time (so the per-step actual-collision probability
$p_{\text{actual}} \ll p_c$, perhaps an order of magnitude smaller because of
the score selection). Then
\begin{equation}\label{eq:sr-bound}
\text{SR}
\;\approx\;
(1 - p_{\text{actual}})^{N_{\text{frame}}}
\;=\;
\exp\!\bigl(-N_{\text{frame}} \cdot p_{\text{actual}}\bigr).
\end{equation}
For $\text{SR} = 0.85$, we need
$p_{\text{actual}} \le -\ln(0.85)/100 \approx 1.63 \times 10^{-3}$.
For $\text{SR} = 0.90$: $p_{\text{actual}} \le 1.05 \times 10^{-3}$.

If $p_{\text{actual}} \approx p_c/k$ for some score-selection improvement
factor $k$, then we need $k \ge p_c / (1.63\times 10^{-3}) = 0.26 / 1.63\times 10^{-3}
\approx 160$.

\subsection{The score-head improvement factor $k$}

This is unknown without closed-loop testing. Empirically YOPO's score head
achieves $k$ on the order of $50$--$200$ in static scenes (the baseline
reports $\sim$$90\%$ static success); whether it carries over to dynamic
scenes is exactly what stage-5 must measure.

\textbf{So the answer to ``is dyn\_dyn = 0.23 enough for $\text{SR} \geq 85\%$?''
is: \emph{depends on $k$}, which we cannot measure without stage-5.}

\section{Decision: stage-5 first does NOT block supervision improvements}

The user's question was: ``if I go straight to stage-5, does it block the
$<10\ms$ and $\geq 85\%$ goals?''

\textbf{It does not block them; it is in fact the only way to measure them.}
The reasoning:

\begin{enumerate}
  \item \textbf{$<10\ms$ goal} depends on inference graph + quantisation +
        device choice (Jetson NX vs Orin Nano). Stage-5 \emph{is} this work.
        No supervision change moves this number; an Option B head with $N=5$
        waypoints would add at most $\sim$$0.2\ms$ to the head (negligible).
        Going stage-5 first decides whether stateless K=10 fits the budget
        without supervision changes. If yes, we ship. If no, we either go to
        stateful (no retrain) or INT8 (no retrain) before pursuing supervision
        upgrades.

  \item \textbf{$\geq 85\%$ SR goal} requires closed-loop sim infrastructure
        that does not yet exist. Building that is part of stage-5. Once it
        exists, we can run the current stage-3.1 model and measure $k$ and
        actual $\text{SR}$. \emph{If} we measure $\text{SR} \geq 85\%$, no
        supervision change is needed --- ship and write the paper.
        \emph{If} we measure $\text{SR} < 85\%$, the gap tells us how
        aggressive the supervision change needs to be (multi-waypoint vs.\
        ESDF vs.\ camera-aware data).

  \item \textbf{Supervision upgrades without a closed-loop bench are
        speculative.} Lowering \texttt{dyn\_dyn} from $0.23$ to $0.16$ via
        Option B is only worthwhile if $\text{SR}$ is actually below $85\%$
        and the bottleneck is the planner's collision-avoidance quality, not
        the controller or the FOV-coverage issue. Without measurement, we
        could spend two weeks on Option B and discover the $k$ factor was
        already strong enough at $\texttt{dyn\_dyn} = 0.23$.
\end{enumerate}

\subsection{Recommended ordering}

\begin{enumerate}
  \item \textbf{Stage-5.A:} ONNX/TensorRT export of the current stage-3.1
        model + Orin NX latency profile. $\approx$3 days.
  \item \textbf{Stage-5.B:} Closed-loop sim harness using \texttt{Simulator/}'s
        existing ROS controller. Run 100 dynamic scenarios; measure $\text{SR}$.
        $\approx$4 days.
  \item \textbf{Decision point.} If $\text{SR} \geq 85\%$: write paper, ship.
        If $\text{SR} < 85\%$: pursue the supervision upgrade with the
        biggest expected return per
        \texttt{01\_collision\_loss\_saturation.tex}'s analysis (likely
        Option B, possibly camera-aware bake).
\end{enumerate}

This sequence \emph{maximises information per week of work}: stage-5 alone
either delivers the paper or tells us exactly which supervision change to
chase. Going supervision-first does the opposite --- it polishes a metric
whose translation to the paper-target SR is uncertain.

\end{document}
